% =====================================================================
%  Lecture 2026-10-01
%
%  Reordered: the proof of the ell_p-ell_q duality theorem follows its
%  statement directly, and the reflexivity block comes after it; in the
%  source the reflexivity block sits between statement and proof. Within
%  that block, the remark introducing the natural map j precedes the
%  definition of reflexivity, which uses j.
%
%  CORRECTED SOURCE ERRORS
%   [p1] finite-dimensional duality proof: gamma_k = f(alpha_k)  ->
%        gamma_k = f(e_k)
%   [p1] identification remark: a norm-preserving bijection  ->  a
%        linear norm-preserving bijection
%   [p2] reflexivity defined by ((X)')' = X  ->  the natural map j is
%        bijective, the source's own reading of that equation
%   [p2] j called the natural injection  ->  the natural map; its
%        injectivity rests on the announced norm equality
%   [p2] bounded variation: sum |w(t_j - t_{j-1})|  ->
%        sum |w(t_j) - w(t_{j-1})|
%   [p2] representation of (C[a,b])': "if w is differentiable"  ->
%        w in C^1[a,b], so that w' is integrable
%   [p2] L_p-L_q duality "for function spaces"  ->
%        (L_p[0,1])' = L_q[0,1] for 1 < p < infinity
%
%  ADDITIONS
%   [p1] notation B(X) = B(X,X), used by the source without definition
%   [p1] finite dimensions: T^{-1} linear as the inverse of a linear
%        bijection; boundedness cited from Example 8.4
%   [p1] Neumann series: B(X) Banach cited from Theorem 9.5; norm
%        convergence S_N -> S turned into pointwise convergence; the
%        telescoping identity written for both orders; continuity of
%        I - T, from Theorem 9.1, used to pass it inside the limit
%   [p1] bounded below: T^{-1} linear
%   [p1] finite-dimensional duality: the standard inner product written
%        out; the map gamma -> f_gamma linear and injective; the case
%        gamma = 0; the maximizer x = gamma / ||gamma||_2
%   [p2] ell_p-ell_q proof: absolute convergence named; the extremal
%        vector replacing the epsilon step; the proof that
%        (f(e_k)) lies in ell_q; the passage f(sum) = sum f via
%        Example 7.5 and continuity of f; linearity and injectivity of
%        beta -> f_beta
%   [p2] natural map: g_x linear, and j linear
%   [p2] reflexive implies Banach: proof supplied, resting on the
%        announced norm equality ||g_x|| = ||x||
%
%  REPORTED BUT UNMODIFIED GAPS
%   [p1] equivalence of injectivity, surjectivity and det T != 0 is
%        left to linear algebra, as in the source
%   [p2] (ell_1)' = ell_infty and (ell_infty)' != ell_1 are stated
%        without proof
%   [p2] ||g_x|| = ||x|| is announced, not proved
%   [p2] reflexivity of ell_p for 1 < p < infinity, and non-reflexivity
%        of ell_1 and ell_infty, are asserted
%   [p2] (L_p[0,1])' = L_q[0,1] is asserted
%   [p2] the Riemann-Stieltjes integral is used without definition;
%        the norm ||w|| is defined and not used; ||f|| is not related
%        to w
%
%  OMITTED
%   [p1] course logistics in the right margin; two quotations
%   [p1] the f <-> gamma box diagram, which the formula carries
%   [p1, p2] doodles unconnected to the material
% =====================================================================

% [p1]
\section{Invertible operators}

\begin{notation}\label{not:BX}
$\Bopone{X} = \Bop{X}{X}$.
\end{notation}

\begin{definition}[Invertible operator, general case]\label{def:invertible-general}
Extending Definition~\ref{def:invertible}, an operator $T \in \Bop{X}{Y}$ is
\emph{invertible} if $T$ is a bijection from $X$ to $Y$ and
$T^{-1} \in \Bop{Y}{X}$.
\end{definition}

Three criteria for invertibility follow.

\begin{proposition}[Finite dimensions]\label{prop:inv-finite}
For a linear map $T \colon \K^n \to \K^n$, that is
$T \in \operatorname{End}(\K^n)$, the following are equivalent.
\begin{enumerate}
  \item $T$ is injective.
  \item $T$ is surjective.
  \item $\det T \ne 0$.
  \item $T$ is invertible.
\end{enumerate}
\end{proposition}

\begin{proof}
The equivalence of (i), (ii) and (iii) is a result of linear algebra. An
invertible $T$ is bijective, so (iv) implies (i). Conversely, if (i), and
hence (ii), holds, then $T$ is a bijection, and $T^{-1}$ is linear as the
inverse of a linear bijection. It is bounded automatically, because $\K^n$
is finite-dimensional, by Example~\ref{ex:finite-dim-bounded}. So $T$ is
invertible.
\end{proof}

\begin{proposition}[Neumann series]\label{prop:neumann}
Let $X$ be a Banach space and $T \in \Bopone{X}$ with $\norm{T} < 1$. Then
$I - T$ is invertible, and
\[
  (I - T)^{-1} = \sum_{k=0}^{\infty} T^k .
\]
\end{proposition}

\begin{remark}
Many calculations with bounded linear operators are similar to those with
complex numbers. For example, if $z \in \C$ and $\abs{z} < 1$, then
$\frac{1}{1-z} = \sum_{k=0}^{\infty} z^k$.
\end{remark}

\begin{proof}[Proof of Proposition~\ref{prop:neumann}]
Let $S_N = \sum_{k=0}^{N} T^k$. For $N > M$,
\[
  \norm{S_N - S_M}
  = \Bigl\lVert \sum_{k=M+1}^{N} T^k \Bigr\rVert
  \le \sum_{k=M+1}^{N} \norm{T^k}
  \le \sum_{k=M+1}^{N} \norm{T}^k ,
\]
which is a piece of the tail of the series $\sum_{k=0}^{\infty} \norm{T}^k$.
That series converges because $\norm{T} < 1$, so given $\eps > 0$ the right
side is less than $\eps$ for all $N > M$ large enough. Thus $(S_N)$ is
Cauchy in $\Bopone{X}$. As $X$ is a Banach space, so is $\Bopone{X}$, by
Theorem~\ref{thm:BXY-banach}, and $S_N$ converges; let
$S = \lim_{N \to \infty} S_N \in \Bopone{X}$. Norm convergence gives
pointwise convergence, since
$\norm{S_N y - S y} \le \norm{S_N - S}\,\norm{y}$ for every $y \in X$.

The sum telescopes, and the two orders of multiplication give the same
result, $S_N(I - T) = (I - T)S_N = I - T^{N+1}$. For $x \in X$,
$\norm{T^{N+1} x} \le \norm{T}^{N+1}\norm{x} \to 0$, so
\[
  S(I-T)x = \lim_{N \to \infty} S_N (I-T) x
  = \lim_{N \to \infty} (I - T^{N+1}) x = x .
\]
The operator $I - T$ is bounded, hence continuous by
Theorem~\ref{thm:bounded-continuous}, and may therefore be taken inside the
limit:
\[
  (I-T)Sx = (I-T)\lim_{N \to \infty} S_N x
  = \lim_{N \to \infty} (I-T)S_N x
  = \lim_{N \to \infty} (I - T^{N+1}) x = x .
\]
So $S(I-T) = (I-T)S = I$. Hence $I - T$ is a bijection with inverse
$S \in \Bopone{X}$.
\end{proof}

\begin{proposition}\label{prop:bounded-below}
Let $T \in \Bop{X}{Y}$ be surjective, and suppose there is a $b > 0$ such
that $\norm{Tx} \ge b\norm{x}$ for all $x \in X$, that is, $T$ is bounded
away from zero. Then $T$ is invertible, and $\norm{T^{-1}} \le \frac{1}{b}$.
\end{proposition}

\begin{proof}
$T$ is injective, that is $\ker T = \set{0}$: if $Tx = 0$, then
$b\norm{x} \le \norm{Tx} = 0$, so $x = 0$. Being also surjective, $T$ is a
bijection, so $T^{-1}$ exists, and it is linear as the inverse of a linear
bijection. It remains to show $T^{-1} \in \Bop{Y}{X}$. For every $y \in Y$
there is an $x \in X$ with $Tx = y$, and then
\[
  \norm{T^{-1} y} = \norm{x} \le \frac{1}{b}\norm{Tx} = \frac{1}{b}\norm{y},
\]
so $\norm{T^{-1}} \le \frac{1}{b}$.
\end{proof}

% [p1]
\section{Dual spaces}

The dual space $X' = \Bop{X}{\K}$ of a normed space $X$ was introduced in
Notation~\ref{not:dual}. In some cases $X'$ can be described explicitly.

\begin{remark}\label{rem:identify}
An equation $(X, \norm{\cdot}) = (Y, \norm{\cdot}')$ between normed spaces
means that the two can be identified with each other: there is a natural
linear bijection $i \colon X \to Y$ with $\norm{i(x)}' = \norm{x}$ for all
$x \in X$.
\end{remark}

\begin{theorem}[Finite dimensions]\label{thm:dual-Rn}
$(\R^n, \norm{\cdot}_2)' = (\R^n, \norm{\cdot}_2)$.
\end{theorem}

\begin{proof}
The identification is given by the standard inner product
$\langle x, \gamma \rangle = \sum_{k=1}^{n} x_k \gamma_k$. Let $f$ be a
linear functional on $(\R^n, \norm{\cdot}_2)$ and $e_1, \dots, e_n$ the
standard basis. For $x = \sum_{k=1}^{n} \alpha_k e_k$ linearity gives
$f(x) = \sum_{k=1}^{n} \alpha_k f(e_k)$. Let $\gamma_k = f(e_k)$; then
\[
  f(x) = \sum_{k=1}^{n} \alpha_k \gamma_k = \langle \alpha, \gamma \rangle ,
\]
where $\alpha = x$, since the coordinates of $x$ in the standard basis are
its entries. So every linear functional is represented by an inner product
with a vector $\gamma \in \R^n$. Conversely, each $\gamma$ defines the
functional $f_\gamma = \langle \cdot\,, \gamma \rangle$, which is bounded by
Example~\ref{ex:finite-dim-bounded}. The map $\gamma \mapsto f_\gamma$ is
linear, and injective because $f_\gamma(e_k) = \gamma_k$. It is therefore a
linear bijection between $\R^n$ and $(\R^n, \norm{\cdot}_2)'$.

It remains to check $\norm{f} = \norm{\gamma}_2$. By the Cauchy--Schwarz
inequality,
\[
  \norm{f} = \sup_{\norm{x}_2 \le 1} \abs{f(x)}
  = \sup_{\norm{x}_2 \le 1} \abs{\langle x, \gamma \rangle}
  \le \sup_{\norm{x}_2 \le 1} \norm{x}_2 \norm{\gamma}_2
  \le \norm{\gamma}_2 .
\]
If $\gamma = 0$ both sides vanish. Otherwise Cauchy--Schwarz attains
equality at $x = \gamma / \norm{\gamma}_2$, which has $\norm{x}_2 = 1$ and
$\langle x, \gamma \rangle = \norm{\gamma}_2$. So
$\norm{f} \ge \norm{\gamma}_2$, and hence $\norm{f} = \norm{\gamma}_2$.
\end{proof}

% [p2]
\begin{theorem}[$\ell_p$--$\ell_q$ duality]\label{thm:lp-duality}
Let $1 < p < \infty$ and $\frac{1}{p} + \frac{1}{q} = 1$. Then
$(\lp{p})' = \lp{q}$. Moreover $(\lp{1})' = \lp{\infty}$, but
$(\lp{\infty})' \ne \lp{1}$.
\end{theorem}

\begin{proof}[Proof of the first statement]
Let $\beta = (\beta_1, \beta_2, \dots) \in \lp{q}$, and define $f_\beta$ on
$\lp{p}$ by
\[
  x = (\alpha_1, \alpha_2, \dots) \longmapsto
  f_\beta(x) = \alpha_1 \beta_1 + \alpha_2 \beta_2 + \dots
  = \sum_{k=1}^{\infty} \alpha_k \beta_k .
\]
By Hölder's inequality, Theorem~\ref{thm:holder-inf},
$\sum_k \abs{\alpha_k \beta_k} \le \norm{x}_p \norm{\beta}_q < \infty$, so
the series converges absolutely. Thus $f_\beta$ is well defined, it is
linear, and
\[
  \norm{f_\beta} = \sup_{\norm{x}_p \le 1} \abs{f_\beta(x)}
  = \sup_{\norm{x}_p \le 1} \Bigl\lvert \sum_k \alpha_k \beta_k \Bigr\rvert
  \le \sup_{\norm{x}_p \le 1} \norm{x}_p \norm{\beta}_q
  \le \norm{\beta}_q .
\]

Equality can be attained in Hölder's inequality. If $\beta = 0$ there is
nothing to show. Otherwise, for each $k$ choose $\sigma_k \in \K$ with
$\abs{\sigma_k} = 1$ and $\sigma_k \beta_k = \abs{\beta_k}$, and put
\[
  \alpha_k = \frac{\sigma_k \abs{\beta_k}^{q-1}}{\norm{\beta}_q^{q-1}} .
\]
From $\frac{1}{p} + \frac{1}{q} = 1$ one gets $(q-1)p = q$, so
$\abs{\alpha_k}^p = \abs{\beta_k}^q / \norm{\beta}_q^q$ and
$\norm{x}_p = 1$, while
\[
  \sum_k \alpha_k \beta_k
  = \frac{\sum_k \abs{\beta_k}^q}{\norm{\beta}_q^{q-1}}
  = \norm{\beta}_q .
\]
Hence $\norm{f_\beta} \ge \norm{\beta}_q$, and $\norm{f_\beta} = \norm{\beta}_q$.

It remains to check whether every $f \in (\lp{p})'$ arises as $f = f_\beta$.
Let $f \in (\lp{p})'$, let $e_k$ be the standard basis vectors, and let
$\beta_k = f(e_k)$. First, $\beta \in \lp{q}$. Fix $N$, choose $\sigma_k$ as
above, and let $x^{(N)}$ have coordinates $\sigma_k \abs{\beta_k}^{q-1}$ for
$k \le N$ and $0$ for $k > N$. By linearity of $f$ on this finite sum, and
by $(q-1)p = q$,
\[
  \sum_{k=1}^{N} \abs{\beta_k}^q = f\bigl(x^{(N)}\bigr)
  \le \norm{f}\,\norm{x^{(N)}}_p
  = \norm{f} \Bigl(\sum_{k=1}^{N} \abs{\beta_k}^q\Bigr)^{1/p} .
\]
If the left side is positive, dividing by its $\frac{1}{p}$-th power gives
$(\sum_{k=1}^{N} \abs{\beta_k}^q)^{1/q} \le \norm{f}$, and the bound holds
trivially otherwise. As $N$ is arbitrary, $\beta \in \lp{q}$.

Second, $f = f_\beta$. Let $x = (\alpha_1, \alpha_2, \dots) \in \lp{p}$. By
Example~\ref{ex:E-schauder}, $x = \sum_k \alpha_k e_k$ with convergence in
$\lp{p}$, and $f$ is continuous by Theorem~\ref{thm:bounded-continuous}, so
\[
  f(x) = f\Bigl(\sum_{k} \alpha_k e_k\Bigr)
  = \lim_{N \to \infty} \sum_{k=1}^{N} \alpha_k f(e_k)
  = \sum_{k} \alpha_k \beta_k = f_\beta(x) .
\]

The map $\beta \mapsto f_\beta$ is linear and, by the norm identity,
norm-preserving, hence injective. It is surjective by the last step. So it is
a linear norm-preserving bijection from $\lp{q}$ onto $(\lp{p})'$, and the
two spaces are identified in the sense of Remark~\ref{rem:identify}.
\end{proof}

\begin{remark}\label{rem:Lp-duality}
The analogous duality holds for function spaces: for $1 < p < \infty$ and
$\frac{1}{p} + \frac{1}{q} = 1$, $(\Lp{p}[0,1])' = \Lp{q}[0,1]$.
\end{remark}

\subsection{Reflexivity}

\begin{remark}[The natural map into the bidual]\label{rem:natural-map}
For $x \in X$ define $g_x \colon X' \to \K$ by $g_x(f) = f(x)$. This is
linear, and $\abs{g_x(f)} = \abs{f(x)} \le \norm{f}\,\norm{x}$ by
Remark~\ref{rem:op-norm-bound}, so $g_x \in X'' = (X')'$ with
\[
  \norm{g_x} = \sup_{\norm{f} \le 1} \abs{g_x(f)}
  = \sup_{\norm{f} \le 1} \abs{f(x)} \le \norm{x} .
\]
The \emph{natural map} $j \colon X \to X''$, $x \mapsto g_x$, is linear. It
will be shown that $\norm{g_x} = \norm{x}$ holds with equality, so that $j$
is norm-preserving, and in particular injective.
\end{remark}

\begin{nfigure}{The natural map $j$ sends $x \in X$ to the functional
$g_x \in X''$, which evaluates every $f \in X'$ at $x$: $g_x(f) = f(x)$.}
\begin{tikzpicture}[scale=0.9]
  \draw[rulegrey, rounded corners=6pt] (0,0) rectangle (1.6,2.4);
  \draw[rulegrey, rounded corners=6pt] (4.4,0) rectangle (6.0,2.4);
  \node at (0.8,2.75) {$X$};
  \node at (5.2,2.75) {$X''$};
  \fill[accent] (0.8,1.2) circle (0.04);
  \node[below] at (0.8,1.15) {$x$};
  \fill[thmaccent] (5.2,1.2) circle (0.04);
  \node[below] at (5.2,1.15) {$g_x$};
  \draw[-{Stealth[length=2mm]}] (0.95,1.25) to[bend left=20]
    node[above] {$j$} (5.05,1.25);
\end{tikzpicture}
\end{nfigure}

\begin{definition}[Reflexive space]\label{def:reflexive}
A normed space $(X, \norm{\cdot})$ is \emph{reflexive} if the natural map
$j \colon X \to X''$ is bijective, that is, if every element of $X''$ can be
written as $g_x$ for some $x \in X$. In the sense of
Remark~\ref{rem:identify} this is written
$((X, \norm{\cdot})')' = (X, \norm{\cdot})$.
\end{definition}

\begin{proposition}\label{prop:reflexive-banach}
Every reflexive space is a Banach space.
\end{proposition}

\begin{proof}
The proof rests on the norm equality $\norm{g_x} = \norm{x}$ announced in
Remark~\ref{rem:natural-map}. The bidual $X'' = \Bop{X'}{\K}$ is a Banach
space by Corollary~\ref{cor:dual-banach}, applied to the normed space $X'$.
Let $(x_n)$ be a Cauchy sequence in $X$. Since
$\norm{g_{x_n} - g_{x_m}} = \norm{j(x_n - x_m)} = \norm{x_n - x_m}$, the
sequence $(g_{x_n})$ is Cauchy in $X''$ and converges to some
$\varphi \in X''$. As $X$ is reflexive, $\varphi = g_x$ for some $x \in X$,
and then $\norm{x_n - x} = \norm{g_{x_n} - g_x} \to 0$. So $X$ is complete.
\end{proof}

\begin{example}\label{ex:reflexive}
$\lp{p}$ is reflexive for $1 < p < \infty$; $\lp{1}$ and $\lp{\infty}$ are
not.
\end{example}

% [p2]
\subsection{The dual of \texorpdfstring{$C[a,b]$}{C[a,b]}}

\begin{definition}[Bounded variation]\label{def:bv}
A function $w \colon [a,b] \to \K$ is \emph{of bounded variation} if
\[
  \operatorname{Var}(w) = \sup \sum_{j=1}^{n} \abs{w(t_j) - w(t_{j-1})}
  < \infty ,
\]
where the supremum is taken over all partitions
$a = t_0 < t_1 < \dots < t_n = b$ of $[a,b]$. Such functions carry the norm
$\norm{w} = \abs{w(a)} + \operatorname{Var}(w)$.
\end{definition}

\begin{theorem}[Representation of {$(C[a,b])'$}]\label{thm:dual-C}
For every $f \in (\Cont{a,b})'$ there is a $w$ of bounded variation on
$[a,b]$ such that
\[
  f(g) = \int_a^b g(t) \dd w(t) \qquad \text{for all } g \in \Cont{a,b} .
\]
The function $w$ is not unique. If $w \in \Contk{1}{a,b}$, then
$f(g) = \int_a^b g(t)\, w'(t) \dd t$.
\end{theorem}
